% Exposé II.B: Erlangen. Sources: theory/framework.md (II.B), site/sections/04-erlangen.html, theory/propositions.json.
\subsection{Erlangen: invariants of group-type methods}\label{sec:erlangen}

\begin{flushright}
\small\itshape A geometry is what a group leaves unchanged. Before asking how a method learns,\\
ask what it cannot help keeping.\\[2pt]
{\upshape\footnotesize after Klein}
\end{flushright}

A method that moves the pretrained weight by a group element can never change that group's invariants. Orbits also
place HRA as a meet, DoRA as a saturation and the scaling methods on a ladder.

\paragraph{In brief}
\begin{itemize}
  \item A method whose weights are always $W_0R$, with $R$ exactly orthogonal, can never change $WW\T$, so neuron norms,
    angles and singular values survive any number of merges in exact arithmetic. HF PEFT's default OFT, BOFT's scale and
    GOFT's default scaler break this.
  \item For injective $W_0$ and even $r\le n-2$, HRA reaches exactly the rank-$r$ LoRA updates of the form $W_0R$,
    $R\in\SO(n)$. Its HF PEFT default initialisation is an apex.
  \item DoRA reaches exactly the weights of LoRA plus a free gain per output neuron. From the same start, any gap between
    the two comes from training.
  \item (IA)$^3$, SSF's scales and DoRA's magnitude lie in the image of rank-one HiRA, $W_0\odot(J+BA)$, which for $W_0$
    without zeros is LoRA seen through $X\mapsto W_0\odot X$.
\end{itemize}

Klein's Erlangen programme of 1872 \citep{xklein1872vergleichende} (\bgref{D.17}) defines a geometry by its group of
motions and studies what the motions leave unchanged. OFT, GOFT, HRA and ETHER multiply the pretrained matrix by an
orthogonal matrix, (IA)$^3$ and SSF's scale multiply it by a diagonal one, and DoRA's magnitude does the same to a LoRA
weight. These are methods of \emph{action type} (\thmref{Definition}{I.2}), $\rho(q)=\gamma(q)\cdot W_0$ with $\gamma$
landing in a group. Whatever that group cannot change, the method cannot change, in training or after any number of
merges.

Klein's programme is the model. \citet{bronstein2021geometric} made it the organising principle of geometric deep
learning, for architectures. Here it is applied to the fine-tuning map.

\begin{analogy*}{what ``Erlangen'' means here}
For a group-type method the Erlangen reading is a theorem, since its image lies in one orbit. LoRA's cone $W_0+\Mr$ has a
singular vertex at $W_0$ (\thmref{Theorem}{II.2}(c)), so it is no orbit; for cones the word names the shape of the image
and its covariance group (\thmref{Definition}{II.14}).
\end{analogy*}

One convention matters. A linear box computes $y=Wx$ with $W\in\R^{m\times n}$, and row $i$ of $W$ is output neuron $i$.
The two Gram matrices, over neurons and over inputs, are
\begin{equation*}
  K_{\mathrm{out}}(W)=WW\T,\qquad K_{\mathrm{in}}(W)=W\T W.
\end{equation*}
A right multiplication $W\mapsto WR$ acts on the input side and fixes $K_{\mathrm{out}}$; a left multiplication
$W\mapsto PW$ acts on the output side and fixes $K_{\mathrm{in}}$. Papers that store $W\T$ swap these words, and we
translate every method into this convention.

% ------------------------------------------------------------------------------------------------
\subsubsection{Complete invariants}\label{sec:erlangen-invariants}

A complete invariant of a group action (\bgref{D.15}) agrees on two matrices exactly when one can be carried to the
other. For the five actions in PEFT it is classical. We write $T_m$ for the invertible diagonal matrices,
$T_m^+=\{\diag(d):d_i>0\}$ for its identity component (\bgref{D.4}), and $\odot$ for the entrywise product.

\begin{theorem}{II.5}{complete invariants of the basic actions; classical}
Let $W,W'\in\R^{m\times n}$.
\begin{enumerate}[label=(\arabic*)]
  \item \emph{Right orthogonal (input side).} $W'=WR$ for some $R\in\Ort(n)$ iff $WW\T=W'W'^{\top}$. For $\SO(n)$, when
    $\rk W=n$, also require that the then unique $R$ have $\det R=1$.
  \item \emph{Left orthogonal.} $W'=PW$ for some $P\in\Ort(m)$ iff $W\T W=W'^{\top}W'$.
  \item \emph{Two-sided.} $W'=PWR$ for some $P\in\Ort(m)$ and $R\in\Ort(n)$ iff $W$ and $W'$ have the same singular
    values, counted with multiplicity.
  \item \emph{Left torus.} If $W$ has no zero row, $W'=DW$ for some $D\in T_m^+$ iff $w_i'\in\R_{>0}\,w_i$ for every $i$.
    For the full torus $T_m$, whose signs are unconstrained, the complete invariant is the set of zero rows together with
    the row \emph{lines}.
  \item \emph{Full Hadamard torus.} $W'=E\odot W$ for some $E\in(\R^\times)^{m\times n}$ iff $W$ and $W'$ have the same
    zero pattern.
\end{enumerate}
\end{theorem}

Item (1) is the orbit-separation consequence of the first fundamental theorem of invariant theory for $\Ort(n)$
(\bgref{D.16}); for a compact group, invariants separate orbits. Equivalently, it is Witt-type extension of isometries.
The two-sided diagonal torus $\{\diag(a)\,X\,\diag(b)\}\cong(T_m\times T_n)/\R^\times$ that appears as LoHa's gauge
(\thmref{Proposition}{II.12}) is a different group from the full Hadamard torus in (5).

\paragraph{Caveat for methods}
No PEFT method realises (4) or (5) as a group action on all of its parameters. (IA)$^3$ entries and DoRA magnitudes are
unconstrained reals initialised at $1$ or at the row norms, so the relevant group in (4) is $T_m$ or its closure.
HiRA's factor $J+BA$ can vanish entrywise. These methods can create new zeros, so for them only a one-way statement
transfers: zeros of $W_0$ remain zero. Table~\ref{tab:erlangen-invariants} lists the five actions with the methods that
live in each.

\begin{table}[t]
\centering\small
\caption{The five basic actions, their complete invariants (\thmref{Theorem}{II.5}) and the methods whose images lie in
one orbit or orbit closure.}
\label{tab:erlangen-invariants}
\begin{tabularx}{\linewidth}{@{}l l >{\raggedright\arraybackslash}p{3.3cm} >{\raggedright\arraybackslash}X@{}}
\toprule
\headrow \textbf{action} & \textbf{group} & \textbf{complete invariant} & \textbf{methods that live here}\\
\midrule
$W\mapsto WR$, input side & $\Ort(n)$ & $K_{\mathrm{out}}=WW\T$ & exact-Cayley OFT and block OFT, HRA, and GOFT and ETHER
  as published\\
$W\mapsto PW$, output side & $\Ort(m)$ & $K_{\mathrm{in}}=W\T W$ & LyCORIS diag-OFT; HF PEFT OFT up to version 0.13\\
$W\mapsto PWR$ & $\Ort(m)\times\Ort(n)$ & singular values with multiplicity & any composite of the two rows above\\
$W\mapsto DW$ & $T_m^+$, $T_m$ & row rays; zero rows and row lines & (IA)$^3$ on $k,v$, SSF's scale, DoRA's magnitude
  (closure only)\\
$W\mapsto E\odot W$ & $(\R^\times)^{m\times n}$ & zero pattern & HiRA (one way only: $Z(W_0)\subseteq Z(W)$)\\
\bottomrule
\end{tabularx}
\end{table}

\begin{explains}
Moving a method from the input side to the output side trades $K_{\mathrm{out}}$ for $K_{\mathrm{in}}$. Letting a scale
change sign trades rays for lines, and letting it reach zero leaves only the zeros of $W_0$.
\end{explains}

The interactive companion has a \emph{Gram lab} for these invariants. It fixes an invertible $4\times4$ matrix $W_0$ with
one zero entry and distinct singular values. The reader picks one of ten methods and drags, types or randomises its
parameters. Each lamp is a float64 test, to $10^{-10}$, that an invariant of $W=\rho(q)$ (the neuron Gram matrix, the
cosines, the input Gram matrix, the singular values, the row lines, the zeros of $W_0$) still equals its value at $W_0$,
and a further lamp tests $\rk\Delta W\le2$. A lamp guaranteed by a theorem is marked as such, and a red lamp names the
hypothesis that was broken. Heatmaps, neuron arrows and strips show $K_{\mathrm{out}}$, the cosines, $K_{\mathrm{in}}$,
$W$, $\sigma(W)$, the rank of $\Delta W$ and $\dim S_1$ against $d$. OFT's Cayley--Neumann switch (HF's default), BOFT's
$s_1\mapsto-s_1$, (IA)$^3$'s $\ell_2\mapsto0$ and DoRA's $B\mapsto0$ each break or restore a hypothesis.

% ------------------------------------------------------------------------------------------------
\subsubsection{What exactly orthogonal methods can never do}\label{sec:erlangen-orthogonal}

If every weight a method reaches is $W_0R$ with $R$ orthogonal, \thmref{Theorem}{II.5}(1) and (3) freeze the neuron Gram
matrix and the singular values, and merging cannot help, since products of orthogonal matrices are orthogonal. The work
lies in saying which implementations are exactly orthogonal, and on which side.

\begin{corollary}{II.6}{what exactly orthogonal methods can never do}
Let a method have image in $W_0\cdot\Ort(n)$ (input side), and work in exact arithmetic. This covers:
\begin{itemize}
  \item OFT and block OFT with an exactly orthogonal parametrisation: the original paper's exact Cayley map, and HF PEFT
    from 0.16 with \texttt{use\_cayley\_\allowbreak neumann=False} on Linear and Conv targets;
  \item GOFT with strict Givens rotations and no scaler, as in the paper (\texttt{no\_scaling=True} in the official
    peft-givens code);
  \item HRA, and ETHER as defined in its paper (not ETHER+);
  \item any number of merge-and-restart cycles of these.
\end{itemize}
Then:
\begin{itemize}
  \item it preserves $K_{\mathrm{out}}$, hence every neuron norm and the angle between any two nonzero neurons, and so the
    hyperspherical energy, which depends only on these angles;
  \item it preserves the whole singular spectrum of $W_0$;
  \item block-diagonal OFT with input blocks $J$, with the same partition in every cycle, preserves each partial Gram
    matrix $W_{:,J}W_{:,J}^{\top}$.
\end{itemize}
\end{corollary}

Several methods that look orthogonal fail the hypothesis. Each item below is a separate computation, and
Table~\ref{tab:erlangen-orthogonal} summarises them.

\begin{remark*}{scope of Corollary II.6; these statements do not follow from its hypothesis}
\begin{itemize}
  \item \emph{BOFT} as implemented in HF PEFT, $W=\diag(s)W_0R\T$ with $R$ exactly orthogonal, has
    $K_{\mathrm{out}}(W)=\diag(s)K_{\mathrm{out}}(W_0)\diag(s)$. So $\cos'_{ij}=\operatorname{sign}(s_is_j)\cos_{ij}$
    whenever $s_is_j\neq0$. It preserves $\lvert\cos\rvert$ when every $s_i$ is nonzero, and preserves $\cos_{ij}$
    whenever $s_is_j>0$. Hence it preserves every cosine when all $s_i$ are nonzero and share a sign (for instance
    $s\in T_m^+$, the constraint in the re-scaled OFT of \citet{qiu2023controlling}, Sec.~3.5); when no cosine of $W_0$
    is zero, that condition is also necessary. HF's \texttt{boft\_s} is an unconstrained real initialised at $1$. At the
    default \texttt{boft\_n\_\allowbreak butterfly\_factor=1} the butterfly permutation is the identity, BOFT is block OFT times an
    output scale, and each input block satisfies
    $W_{:,J}W_{:,J}^{\top}=\diag(s)(W_0)_{:,J}(W_0)_{:,J}^{\top}\diag(s)$. HF PEFT OFT 0.14 to 0.15 (exact Cayley, input
    side) carried a trainable output scale \texttt{oft\_s} and has the same form.
  \item \emph{GOFT's official code} by default (\texttt{no\_scaling=False}) adds a trainable input-side diagonal $g$,
    initialised at ones. Then $W=W_0\diag(g)G$ and $K_{\mathrm{out}}(W)=W_0\diag(g)^2W_0\T$, whatever the Givens part $G$
    is. This is the paper's GOFT* ablation, and it is not covered.
  \item \emph{Output-side implementations} (LyCORIS diag-OFT at its defaults; HF PEFT OFT up to 0.13, which merged to
    $R\T W_0$; HF OFT on embedding layers; ETHER with \texttt{flip\_side=True}) have image in $\Ort(m)\cdot W_0$. They
    preserve $K_{\mathrm{in}}$, the partial input Grams $W_{I,:}^{\top}W_{I,:}$ and the spectrum, and they do not
    preserve $K_{\mathrm{out}}$. ETHER's official code with \texttt{ether\_nb>1} applies a different reflection to each
    output-row group, which keeps neuron norms and within-group angles but neither the cross-group cosines nor the
    spectrum.
  \item \emph{HF PEFT's default OFT} since 0.18 (\texttt{use\_cayley\_\allowbreak neumann=True}, 5 terms) builds
    $R=I+2Q+2Q^2+2Q^3+Q^4=C(Q)(I-Q^4)$, where $C(Q)=(I+Q)(I-Q)^{-1}$ is the exact Cayley map (HF's exact branch
    evaluated at $-Q$). Hence $R\T R=(I-Q^4)^2\ne I$, and the singular values of $R$ are $\lvert1-\theta_k^4\rvert$,
    where $\pm i\theta_k$ are the eigenvalues of $Q$: a systematic contraction for small $Q$, and $R$ is singular when
    some $\theta_k=1$. Already the adapted weight has $K_{\mathrm{out}}=W_0(I-Q^4)^2W_0\T$, a drift of order
    $\lVert Q\rVert^4$, and merge cycles compound it. For $\lVert Q\rVert_{\mathrm{op}}\le2^{1/4}$ every factor is a
    contraction, so the drift is a monotone shrinkage of neuron norms and singular values (\thmref{Theorem}{V.3},
    corollary 4). In 0.16 and 0.17 the last coefficient was $2$, so $R\T R=I+4Q^6+4Q^8$, still not orthogonal, with
    drift of order $\lVert Q\rVert^6$. (HF's docstring says the default is \texttt{False}; the dataclass default is
    \texttt{True}.)
  \item \emph{Quasi-orthogonal variants} preserve neither $K_{\mathrm{out}}$ nor its normalised form exactly: qGOFT,
    whose $2\times2$ blocks are arbitrary matrices with only a soft orthogonality penalty; ETHER+, which is two-sided and
    whose factor $I-uu\T+vv\T$ is singular when $u\perp v$; and the default HF OFT above. Once $n>2$, a single
    non-orthogonal $2\times2$ block, conformal or not, generically changes some cosine.
  \item \emph{Finite precision.} The invariance is exact only in exact arithmetic; low-precision merges add a rounding
    drift that grows like a random walk over merge cycles.
\end{itemize}
\end{remark*}

\paragraph{Numerical check}
With $s=(-0.5,1,\dots,1)$, BOFT changes a cosine by $1.5$ while $\lvert\cos\rvert$ is fixed to $4\times10^{-16}$, and an
all-negative $s$ fixes every cosine to $10^{-15}$. GOFT with a scaler $g=1\pm0.05$ moves cosines and singular values. A
conformal block $\sqrt2\,\mathrm{Rot}(0.3)$ on two of $8$ input coordinates changes some cosine, by $0.17$. HF's default
OFT satisfies $R\T R=(I-Q^4)^2$ to $5\times10^{-16}$, with $\lVert R\T R-I\rVert=1.5\times10^{-2}$ already for block
size $8$ and skew entries of size about $0.07$, and it is singular at $\theta=1$ (Appendix~\ref{app:repro}).

\begin{table}[t]
\centering\footnotesize
\caption{Implementations that look orthogonal, read against \thmref{Corollary}{II.6}. Each row is a separate computation
(see the remarks after the corollary).}
\label{tab:erlangen-orthogonal}
\begin{tabularx}{\linewidth}{@{}>{\raggedright\arraybackslash}p{2.9cm} >{\raggedright\arraybackslash}p{3.4cm}
  >{\raggedright\arraybackslash}p{1.7cm} >{\raggedright\arraybackslash}X@{}}
\toprule
\headrow \textbf{implementation} & \textbf{map} & \textbf{exactly orthogonal?} & \textbf{what it keeps}\\
\midrule
BOFT in HF PEFT; HF PEFT OFT 0.14--0.15 & $\diag(s)\,W_0R\T$, $R$ exactly orthogonal, $s$ free, initialised at $1$ & only $R$
  & $K_{\mathrm{out}}\mapsto\diag(s)K_{\mathrm{out}}\diag(s)$, and block Grams likewise. $\lvert\cos_{ij}\rvert$
  survives when $s_is_j\ne0$, $\cos_{ij}$ when $s_is_j>0$. At HF's default single butterfly factor, BOFT is block OFT
  times $s$.\\
GOFT, official code default & $W_0\,\diag(g)\,G$, with a trainable input scale $g$ & no & Nothing on this list, since
  $K_{\mathrm{out}}=W_0\diag(g)^2W_0\T$.\\
output-side OFT (LyCORIS diag-OFT at its defaults; HF PEFT OFT up to 0.13) & image in $\Ort(m)\cdot W_0$ & yes, on the
  other side & $K_{\mathrm{in}}$, the partial input Grams and the spectrum. Not $K_{\mathrm{out}}$.\\
HF PEFT default OFT since 0.18 (\texttt{use\_cayley\_\allowbreak neumann=True}) & $R=I+2Q+2Q^2+2Q^3+Q^4=C(Q)(I-Q^4)$, with $C$ an
  exact Cayley map & no, $R\T R=(I-Q^4)^2$ & Nothing exactly. The adapted weight already has
  $K_{\mathrm{out}}=W_0(I-Q^4)^2W_0\T$, and for small $Q$ each merge shrinks $K_{\mathrm{out}}$ and the spectrum
  further. (In 0.16--0.17, $R\T R=I+4Q^6+4Q^8$.)\\
qGOFT & free $2\times2$ blocks with a soft orthogonality penalty & no & Neither $K_{\mathrm{out}}$ nor the cosines. Once
  $n>2$, one non-orthogonal block, conformal or not, generically moves a cosine.\\
ETHER+ & two-sided, $I-uu\T+vv\T$, singular when $u\perp v$ & no & Neither $K_{\mathrm{out}}$ nor the cosines.\\
\bottomrule
\end{tabularx}
\end{table}

\paragraph{Credit}
\citet{qiu2023controlling} introduced orthogonal fine-tuning, motivated by preserving hyperspherical energy
\citep{liu2018learning}; their re-scaled OFT (Sec.~3.5) adds the positive per-neuron scale that HF BOFT carries. BOFT is
due to \citet{liu2023parameter}, building on butterfly factorisations \citep{dao2019learning}; GOFT and qGOFT are due to
\citet{ma2024parameter}, with an official peft-givens implementation; ETHER and ETHER+ to \citet{bini2024ether}; HRA to
\citet{yuan2024bridging}. LyCORIS is described by \citet{yeh2023navigating}. The Cayley--Neumann default, BOFT's scale
and the output-side merge are read from the Hugging Face PEFT source \citep{xmangrulkar2022peft} (the OFT configuration
and the BOFT layer).

\begin{explains}
``Orthogonal fine-tuning preserves knowledge'' corresponds to preserving a complete invariant: in exact arithmetic an
exactly orthogonal method can \emph{never} change spectra or neuron geometry, through any number of merges
(\thmref{Theorem}{V.3}, corollary 7), up to rounding. So a variant that must change them adds a non-orthogonal factor
(the per-neuron scale of Qiu et al.'s re-scaled OFT, as in HF BOFT and HF OFT 0.14 to 0.15, or GOFT's default input
scaler). HF's default OFT is only close to orthogonal; with block size $8$ and fresh skew generators with entries near $0.05$ in
every cycle, its merges shrink $K_{\mathrm{out}}$ by about $0.2\%$, $1.3\%$, $6\%$ and $22\%$ after $1$, $10$, $50$ and
$200$ cycles (relative Frobenius change, averaged over draws). The statement
concerns reach. It does not show that preserving the invariant causes knowledge retention, and the control study of
ETHER \citep{bini2024ether} found a non-orthogonal multiplicative variant comparably stable. Earlier claims that BOFT or
qGOFT preserve the full Gram matrix are corrected here.
\end{explains}

% ------------------------------------------------------------------------------------------------
\subsubsection{HRA is a meet of two images}\label{sec:erlangen-hra}

A Householder reflection $H_u=I-2uu\T/\lVert u\rVert^2$ is orthogonal and differs from $I$ by rank one. So
$W_0H_{u_1}\cdots H_{u_r}$ lies in the orthogonal orbit of $W_0$ and in LoRA's cone $W_0+\Mr$, and for injective $W_0$
and even $r$ the Cartan--Dieudonn\'e theorem (\bgref{D.11}) says that HRA fills the cone's intersection with the
rotation orbit. HF PEFT repeats each reflection at initialisation, so to first order each pair can only turn its own
vector against the rest of the space, and HRA starts at an apex, as zero-init LoRA does (\thmref{Theorem}{II.2}).

\begin{theorem}{II.7}{HRA is the meet of LoRA with the rotation orbit, and it starts at an apex}
Let $W_0\in\R^{m\times n}$ be injective ($\rk W_0=n$), and let $r$ be even with $2\le r\le n-2$. Define
$\mathrm{HRA}_r:(u_1,\dots,u_r)\mapsto W_0H_{u_1}\cdots H_{u_r}$ with $u_i\ne0$, where
$H_u=I-2uu\T/\lVert u\rVert^2$. Write $\mathrm{OFT}_{\SO}:=W_0\cdot\SO(n)$ for the rotation orbit. (The $\Ort(n)$-orbit
would add the odd component: $W_0H_{e_1}$ lies in $W_0+\Mr$ but not in $\Im\mathrm{HRA}_r$.)
\begin{enumerate}[label=(\alph*)]
  \item \textbf{Meet.} $\Im\mathrm{HRA}_r=(W_0+\Mr)\cap W_0\cdot\SO(n)$. Further, $d(\mathrm{HRA}_r)=rn-\tfrac{r(r+1)}2$,
    so the generic fibre has dimension $\tfrac{r(r+1)}2=rn-d$. This gauge is non-linear. The evident continuous global
    symmetries are the $r$ scalings $u_i\mapsto c_iu_i$; the braid group also acts globally by Hurwitz moves
    $(u_i,u_{i+1})\mapsto(H_{u_i}u_{i+1},u_i)$, since $H_aH_b=H_{H_ab}H_a$.
  \item \textbf{First-order collapse.} At the paired initialisation $u_{2i-1}=u_{2i}=w_i$ ($i\le k=r/2$, with the $w_i$
    independent and $U=\spanop\{w_i\}$),
    \begin{equation*}
      S_1=W_0\cdot\{\Omega\in\so(n):P_{U^\perp}\Omega P_{U^\perp}=0\},\qquad\dim S_1=kn-\tfrac{k(k+1)}2<d(\mathrm{HRA}_r).
    \end{equation*}
  \item \textbf{Vertex.} The germ of $\Im\mathrm{HRA}_r$ at $W_0$ is diffeomorphic to the germ at $0$ of
    $\Sigma_r=\{\Omega\in\so(n):\rk\Omega\le r\}$, the skew-determinantal cone (\bgref{C.13}), which is not linear. So
    $W_0$ is a singular point (the vertex) of the image. Together with (b), the paired initialisation is an \textbf{apex}.
  \item \textbf{Scope.}
    \begin{itemize}
      \item \emph{Not a product.} (a) is a meet in the poset of subsets of $\Theta$ containing $W_0$
        (\thmref{Observation}{I.4}(c)). HRA is not the
        categorical product of $\mathrm{LoRA}_r$ and $\mathrm{OFT}_{\SO}$: the set-level fibre product
        $\{(B,A,R):W_0+BA=W_0R\}$ has local dimension $d+r^2=rn+\tfrac{r(r-1)}2\ne rn$ over a generic point of the
        meet. (It is not pure-dimensional: over $W_0$ alone it contains $\{(B,0,I)\}$, of dimension $mr$.)
      \item \emph{Exactly orthogonal OFT is smaller.} The Cayley map never produces a rotation with eigenvalue $-1$.
        For exact full-block Cayley OFT (the original paper; HF with \texttt{use\_cayley\_\allowbreak neumann=False}),
        $(W_0+\Mr)\cap\Im\mathrm{OFT}$ is dense in $\Im\mathrm{HRA}_r$ but not equal to it: $W_0H_{e_1}H_{e_2}$, with
        $H_{e_1}H_{e_2}=\diag(-1,-1,1,\dots,1)$, is missing. For exact block Cayley OFT with $b<n$ the intersection is a
        proper subset of $\Im\mathrm{HRA}_r$.
      \item \emph{HF's default OFT is incomparable.} HF PEFT's default OFT is not orthogonal (\thmref{Corollary}{II.6};
        \texttt{use\_cayley\_\allowbreak neumann=True}). Since $R\T R=(I-Q^4)^2=I$ forces $Q=0$ near $Q=0$, its image meets
        $\Im\mathrm{HRA}_r$ near $W_0$ only at $W_0$, while its intersection with $W_0+\Mr$ contains non-orthogonal
        points such as $W_0(I+2Q+2Q^2+2Q^3+Q^4)$ with $\rk Q=2$.
      \item \emph{Implementation.} The theorem concerns HRA with $\lambda<\infty$ (HF \texttt{apply\_GS=False}). The
        strictly orthogonal Gram--Schmidt variant has the smaller image $W_0(I-2P_U)$, of dimension $r(n-r)$. The
        paired initialisation is HF's default (\texttt{repeat\_interleave}) and is not stated in the HRA paper. HF HRA
        acts on \texttt{in\_features}, so injectivity requires \texttt{out\_features}${}\ge{}$\texttt{in\_features}
        (and full column rank); the theorem does not cover \texttt{down\_proj} or GQA key/value projections. The
        Gram--Schmidt variant never contains $W_0$, since $I-2P_U\ne I$, and at HF's default paired initialisation it
        runs Gram--Schmidt on duplicated columns, which is numerically degenerate.
    \end{itemize}
\end{enumerate}
\end{theorem}

\paragraph{Numerical check}
At $n=9$, the paired-initialisation Jacobian has rank $8$ against generic rank $15$ for $r=2$, and $15$ against $26$
for $r=4$ (Appendix~\ref{app:repro}).

\paragraph{Credit}
The sharp Cartan--Dieudonn\'e theorem, that a rotation moving a $k$-dimensional space is a product of exactly $k$
reflections, is due to \citet{xscherk1950decomposition}; see \citet{xartin1957geometric}. Householder parametrisations
in deep learning go back to \citet{mhammedi2016efficient}, and Lie-algebra charts such as the Cayley map to
\citet{lezcanocasado2019cheap}. \citet{yuan2024bridging} proposed HRA as a bridge between low-rank and orthogonal
adaptation, which the theorem makes precise. The paired initialisation is read from the HF PEFT HRA layer
\citep{xmangrulkar2022peft}.

\begin{explains}
\begin{itemize}[leftmargin=1.2em]
  \item HRA's ``bridge between low-rank and orthogonal adaptation'' is, at the level of images, the meet of LoRA with the
    rotation orbit $W_0\cdot\SO(n)$. It is a meet of images and no product of methods. It strictly contains
    LoRA${}\cap{}$OFT for exactly orthogonal (Cayley) OFT, as in the original paper or HF with
    \texttt{use\_cayley\_\allowbreak neumann=False}; for HF's default Cayley--Neumann OFT the two are incomparable. Compare
    \thmref{Proposition}{I.8}(a), where LoRA-XS \citep{baazy2024lora} is a genuine categorical product
    $\mathrm{FA}\times\mathrm{FB}$.
  \item HF recommends an even $r$ because its paired initialisation needs pairs, and for odd $r$ it falls back to a
    random Kaiming initialisation. The determinant identity $\det R=(-1)^r$ explains why no odd-$r$ HRA can start at
    $W_0$, so that fallback is a pointing defect.
  \item New here (earlier analyses labelled this base point smooth): \textbf{HRA's default initialisation is an apex}. At
    first order, each pair can only rotate its own vector against the rest of the space. This is the multiplicative twin
    of LoRA's dead $A$-direction (\thmref{Lemma}{II.4}), so HRA belongs in the apex column beside LoRA. Whether the
    collapse slows early training is open (\S\ref{sec:analogy}).
  \item Rebasing commutes with the meet, so after $K$ cycles re-HRA reaches $(W_0+\mathcal M_{\le Kr})\cap W_0\cdot\SO(n)$
    (\thmref{Theorem}{V.3}, corollary 5).
\end{itemize}
\end{explains}

In the Gram lab, choosing HRA and resetting to $q_0$ gives $\dim S_1=3$ against $d=5$, the values of (a) and (b) at
$n=4$, $r=2$. Moving $u_2$ keeps the Gram, spectrum and rank lamps lit together, and the preset $u=(e_1,e_2)$ lands on
$W_0H_{e_1}H_{e_2}$, which exact Cayley OFT never reaches.

% ------------------------------------------------------------------------------------------------
\subsubsection{LOFT fixes one leaf of the meet}\label{sec:erlangen-loft}

The proof of \thmref{Theorem}{II.7} (\S\ref{proof:II.7}) already shows how the meet is built. It is the union of
the orbits $W_0\cdot\SO(V)$ over the $r$-dimensional input subspaces $V$, one subgroup orbit for each point of the
Grassmannian $\Gr(r,n)$ (\bgref{C.15}). Call each orbit a \emph{leaf}. Leaves overlap along lower-rank strata, so the
leaves cover the meet without partitioning it; a generic point lies on exactly one leaf, fixed by its moved space
$\operatorname{im}(R-I)$. HRA's reflection vectors move the leaf as it trains. LOFT \citep{zhao2026loft} freezes the
leaf and makes the choice of leaf the design variable, which the paper calls the adaptation support.

\begin{remark*}{2026: LoCO}
One component of LoCO \citep{nguyen2026loco} is the Cayley rotation of a low-rank skew generator $UV\T-VU\T$, with
$U,V\in\R^{n\times r}$. For injective $W_0$ and $2r\le n-2$, its image is $\Im\mathrm{HRA}_{2r}$ minus the points $W_0R$
at which $R$ has eigenvalue $-1$, so like HRA it moves its leaf while it trains. Writing the generator as $ZJZ\T$ with
$Z=[U\ V]$ shows that its generic fibre is an orbit of the linear group $\mathrm{Sp}(2r,\R)$, acting by $Z\mapsto ZM$.
The non-linear gauge of \thmref{Theorem}{II.7}(a) is therefore a property of the reflection parametrisation; LoCO's
component reaches an open dense part of the same set with the linear gauge $\mathrm{Sp}(2r,\R)$. A sum of several LoCO components is in general only approximately
orthogonal, and then leaves the meet.
\end{remark*}

\begin{example*}{LOFT: a leaf, and how to choose it}
Let $P_r\in\R^{r\times n}$ have orthonormal rows spanning $V$, and let $\Ort(V)=\{I+P_r\T(T-I)P_r:T\in\Ort(r)\}$ be the
orthogonal maps that fix $V^\perp$ pointwise. LOFT trains $\rho(E)=W_0\big(I+P_r\T(Q(E)-I)P_r\big)$ with $E\in\so(r)$,
$E_0=0$, so $\lvert M\rvert=r(r-1)/2$. We read it with an exactly orthogonal $Q$ (the exponential, or the Cayley map
$C$); the released code defaults to a five-term Cayley--Neumann series with $R\T R=(I-X^6)^2$, which is only
approximately orthogonal, so like HF's default OFT (\thmref{Corollary}{II.6}) that version leaves its leaf.
\begin{enumerate}[label=(\alph*)]
  \item \emph{One leaf.} If $W_0P_r\T$ has rank $r$, then
    $\big(W_0+W_0P_r\T\,\R^{r\times r}P_r\big)\cap W_0\cdot\Ort(n)=W_0\cdot\Ort(V)$. For the principal support
    $P_r=V_r\T$ the first set is the image of LoRA-XS (this needs $\sigma_r(W_0)>0$), so the leaf is
    LoRA-XS${}\cap W_0\cdot\Ort(n)$. Principal-support LOFT reaches an open dense part of its component $W_0\cdot\SO(V)$;
    for injective $W_0$ and a Cayley $Q$, its image is exactly LoRA-XS${}\cap{}$the image of exact full-block Cayley
    OFT.
  \item \emph{Leaves over the Grassmannian.} If $W_0$ is injective,
    $(W_0+\Mr)\cap W_0\cdot\SO(n)=\bigcup_{V\in\Gr(r,n)}W_0\cdot\SO(V)$, which for even $r$ is $\Im\mathrm{HRA}_r$
    (\thmref{Theorem}{II.7}(a)). LOFT's image is the whole leaf $W_0\cdot\SO(V)$ with the exponential, and an open dense
    part of it with a Cayley map, which never reaches eigenvalue $-1$.
  \item \emph{A regular base point.} If $W_0P_r\T$ has rank $r$, then $S_1=W_0P_r\T\so(r)P_r$ has dimension
    $r(r-1)/2=d(M)$: LOFT starts at a smooth point of its leaf and has no gauge, while HRA's paired initialisation is an
    apex (\thmref{Theorem}{II.7}(b)). For $r\ge4$ there is no smooth pointed arrow LOFT${}\to{}$HRA at that
    initialisation, although the images are nested (obstruction test O3 of \S\ref{sec:atlas}).
  \item \emph{Choosing the leaf.} If $Q'(0)=\id$ (the exponential, or $E\mapsto C(E/2)$), the gradient pulled back to $E$
    at $\theta_0$ is $P_r\skewop(W_0\T G)P_r\T$ with $G=\nabla L(W_0)$; the Cayley map $C$ itself doubles it, which does
    not change the best support. Its squared norm is at most $2\sum_{k\le r/2}\mu_k^2$, where $\pm i\mu_k$ are the
    eigenvalues of $\skewop(W_0\T G)$, with equality when $V$ spans the invariant subspace of the top $\lfloor r/2\rfloor$
    pairs.
\end{enumerate}
\end{example*}

Appendix~\ref{proof:loft} proves (a)--(d).

\paragraph{Numerical check}
For $n=8$, $m=10$ and $r=2,4,6$, the leaf identity, the union, the arrows into exact Cayley OFT and LoRA-XS, and
$\dim S_1=r(r-1)/2$ all hold to about $10^{-15}$ (Appendix~\ref{app:repro}).

\paragraph{Credit}
LOFT, the SkewGrad and GradSVD supports, part (d) and its bound (their Proposition~2), the reading of each support as
selecting a Lie subgroup (their Remark~1) and the suggestion of dynamic supports (their Remark~3) are from
\citet{zhao2026loft}. Their Table~1 writes full and block OFT, GOFT, BOFT, HRA and PSOFT as choices of support, width,
depth and transform in a multi-factor version of the same primitive. Parts (a)--(c) place the single-factor method inside
\thmref{Theorem}{II.7}. PSOFT is due to \citet{wu2025efficient}; LoRA-XS to \citet{baazy2024lora}.

\begin{explains}
Orthogonal methods differ in which leaves they use and how they combine them. Block OFT, GOFT and BOFT multiply rotations
of fixed coordinate subspaces; HRA moves its leaf while it trains; PSOFT takes the principal leaf and adds two tori; LOFT
takes one leaf, chosen before training from $W_0$ or from the gradient at $\theta_0$. By part (d), SkewGrad is the
orthogonal counterpart of the alignment mechanism of \thmref{Proposition}{II.3}, and GradSVD transplants the
gradient-based frames of LoRA-GA \citep{wang2024lorab}. With the leaf frozen, merge-and-restart gains nothing. With
finitely many supports revisited across restarts, the reach is the orbit of the connected group whose Lie algebra
(\bgref{D.8}) the $\so(V_k)$ generate (\thmref{Theorem}{V.3}(b)); numerically, generic supports with $r\ge2$ generate all
of $\SO(n)$ once $\lceil n/r\rceil$ of them are used. The principal support does not move when it is re-selected after
a merge, and data-chosen supports have to be checked case by case.
\end{explains}

% ------------------------------------------------------------------------------------------------
\subsubsection{DoRA is a torus saturation of LoRA}\label{sec:erlangen-dora}

DoRA normalises each row of LoRA's matrix $V=W_0+BA$ and multiplies row $i$ by a learned magnitude $\mu_i$. That is one
row scaling, by $\mu_i/\lVert v_i\rVert$, so every DoRA weight is a diagonal matrix times a LoRA weight, and conversely
up to a repair for vanishing rows. We write $\mathrm{Diag}_m$ for the monoid of diagonal matrices (\bgref{D.2}),
$\mathrm{Mon}_m=T_m\rtimes S_m$ for the monomial matrices (\bgref{D.3}), and $N$ for row normalisation.

\begin{proposition}{II.8}{DoRA is LoRA followed by an output scale, at the level of images}
Assume $1\le r\le\min(m,n)$ and $s\ne0$. Let $N$ normalise rows; DoRA's magnitude $\mu\in\R^m$ is per output neuron, as
in HF PEFT and the official code (the paper's notation says ``column-wise'', in the transposed convention). Images refer
to the merged (eval-mode) weight; HF's train-mode forward with \texttt{lora\_dropout}${}>0$ is not of the form $Wx$. Let
$W_0$ have no zero row, let $\mathcal U$ be the open set of matrices with no zero row, and let $\mathrm{Diag}_m$ be the
monoid of all diagonal matrices. Then
\begin{equation*}
  \Im\mathrm{DoRA}_r=\{\diag(\mu)\,N(V):V\in(W_0+\Mr)\cap\mathcal U,\ \mu\in\R^m\}=\mathrm{Diag}_m\cdot(W_0+\Mr),
\end{equation*}
which is exactly the image of ``$\mathrm{LoRA}_r$ followed by a trainable output scale $\ell\in\R^m$''. With $\mu$
restricted to $(\R^\times)^m$, or equivalently after intersecting with $\mathcal U$:
$\Im\mathrm{DoRA}_r\cap\mathcal U=T_m\cdot(W_0+\Mr)\cap\mathcal U$. Moreover:
\begin{itemize}
  \item $d(\mathrm{DoRA}_r)\le\min\big(mn,\ r(m+n-r)+m\big)$. At a parameter point with $\mu\in(\R^\times)^m$ and
    $X=sBA\in\MM_r$, with $U=\operatorname{col}X$ and $R=\operatorname{row}X$, the Jacobian rank is
    $r(m+n-r)+m-\dim\ker\big(\delta\mapsto P_{U^\perp}\diag(\delta)W_0P_{R^\perp}\big)$; where some $\mu_i=0$ it is
    smaller when $r\ge2$, and for $r=1$ when some row $w_{0,i}$ with $\mu_i=0$ lies outside $R$. For generic $W_0$, $d=\min\big(mn,\ r(m+n-r)+m\big)$; in particular $d=r(m+n-r)+m$ iff
    $(m-r)(n-r)\ge m$, which holds for all practical shapes. For special $W_0$ it is smaller: for $W_0=ac\T$ with all
    $a_i\ne0$ and $r<n$, $d=r(m+n-r)+m-r$, strictly below the generic value iff $r<m$ and $r\le n-2$.
  \item The image covariance group (\thmref{Definition}{II.14}) contains $\mathrm{Mon}_m\times\GL_n$.
  \item At $B=0$, in HF's magnitude coordinates, $\Delta W=\big(\diag(\mu_i/\lVert w_{0,i}\rVert)-I\big)W_0$, which can
    have rank up to $\rk W_0$, so DoRA escapes LoRA's rank bound.
\end{itemize}
\end{proposition}

\paragraph{Numerical check}
The Jacobian rank is $9=mn$ at $(m,n,r)=(3,3,2)$, where $r(m+n-r)+m$ gives $11$; it is $36$ at $(8,8,2)$, matching, and
$34$ there when one $\mu_i=0$; it is $24$ for a rank-one $W_0$ at $(6,6,2)$ (Appendix~\ref{app:repro}).

\paragraph{Credit}
DoRA is due to \citet{liu2024dora}. Its magnitude--direction split follows weight normalisation
\citep{salimans2016weight}. From a different direction, \citet{zhang2025primacy} argue that the magnitude of the update
is the main driver of LoRA's performance.

\begin{explains}
At the level of images, DoRA's ``magnitude--direction decoupling'' is LoRA plus $m$ free output gains. What DoRA changes
beyond that lives in its training dynamics. Under gradient flow or SGD there are two different geometries. The exact
gradient of the normalised parametrisation projects the direction gradient row by row, which is a change of metric. The
detached backward of the DoRA paper (Sec.~4.3) and of HF PEFT (\texttt{weight\_norm.detach()}) removes that projection;
the resulting update is a non-gradient pseudo-gradient that no metric describes, and it need not be a descent direction.
Both the paper and HF train with AdamW, under which neither variant follows a fixed metric, and the coordinates matter
in their own right: $\mu_i$ starts at $\lVert w_{0,i}\rVert$ and $\ell_i$ at $1$, so an equal step on $\mu_i$ and on a
gain $\ell_i$ changes the row gain by different amounts.
\end{explains}

\paragraph{Experiment D}
With one optimizer (AdamW, same learning rate, weight decay and data order) and the same $W_0$, compare four arms:
(i) LoRA + gain $\ell$ with $\ell_0=\mathbf 1$; (ii) LoRA + gain in DoRA's coordinates,
$W=\diag(\mu_i/\lVert w_{0,i}\rVert)(W_0+sBA)$ with $\mu_0=(\lVert w_{0,i}\rVert)_i$, which isolates the
optimizer-coordinate effect without normalisation; (iii) DoRA with the exact backward; (iv) DoRA with HF's detached
backward. Since the images are equal, any gap between arms comes from the parametrisation, the backward rule or the
optimizer coordinates, and the four arms separate these. The DoRA paper's own ablation reports (iii) and (iv) within
$0.2$ points. \thmref{Prediction}{X.2} uses the equal images to set up this controlled comparison.

The covariance group $\mathrm{Mon}_m\times\GL_n$ decides where \thmref{Observation}{II.15} applies. One orthogonal
rotation of the residual stream of a gain-fused pre-norm model, as in SliceGPT \citep{xashkboos2024slicegpt}, acts on the
input side of $q$, $k$, $v$, gate and up, where DoRA is $\GL_n$-covariant, so DoRA on those boxes reaches the same
functions after the rotation as before; on $o$ and down the sufficient condition fails. Training outcomes are a separate
question.

% ------------------------------------------------------------------------------------------------
\subsubsection{The torus ladder}\label{sec:erlangen-ladder}

Write a weight as $W_0\odot(J+Z)$, with $J$ the all-ones matrix. Output scalings use $Z=(\ell-\mathbf 1)\mathbf 1\T$, of
rank one; two-sided scalings use $ab\T-J$, of rank at most two; HiRA allows any $Z=BA$ of rank at most $r$. These
choices form a ladder of images, and no rung can fill a zero of $W_0$.

\begin{proposition}{II.9}{the torus ladder}
\begin{enumerate}[label=(\alph*)]
  \item \textbf{HiRA.} $\rho(B,A)=W_0\odot(J+BA)$, where $J$ is the all-ones matrix; HF PEFT has $B_0=0$ and no scale
    factor on Linear and Conv layers (on embeddings it uses $A_0=0$). For every $W_0$,
    $\Im\mathrm{HiRA}_r=W_0+D_{W_0}(\Mr)$ with $D_{W_0}:X\mapsto W_0\odot X$. If $W_0$ has no zero entries,
    $D_{W_0}\in\GL(\Theta)$, so inside the slice at $\theta_0$ HiRA is LoRA transported by a diagonal automorphism: it has
    the same $d=r(m+n-r)$ (for $r\le\min(m,n)$) and the same gauge. In all cases $\rk\Delta W\le r\cdot\rk W_0$. The
    transport depends on $\theta_0$ (\thmref{Proposition}{II.13}).
  \item \textbf{Scalings.} Output-side scalings, $\diag(\ell)W_0=W_0\odot(J+(\ell-\mathbf 1)\mathbf 1\T)$, and
    input-side scalings, $W_0\diag(\ell)=W_0\odot(J+\mathbf 1(\ell-\mathbf 1)\T)$, have images in $\Im\mathrm{HiRA}_1$.
    Output-side examples: (IA)$^3$ on keys and values, the post-linear scale of SSF, DoRA's magnitude. Input-side
    examples: (IA)$^3$'s feed-forward vector, which scales the input of $W_{\mathrm{down}}$, and a LayerNorm gain
    $\gamma$ folded into the next linear map. The fold is exact when the norm's output is consumed only by linear maps:
    the input norms of pre-norm blocks, and a final norm before an untied head (in post-norm BERT, the last block's norm
    feeds only the head). It does not apply to Gemma-style post-sublayer norms, which HF's default LN-tuning targets for
    Gemma~2 and~3, to QK-norms, or to post-norm residual paths. Where it applies, $\gamma$ is tied across all consumers of
    the norm (Q, K and V, or gate and up) rather than being a free point per box, and with tied embeddings a fold into
    the head breaks the tie. LayerNorm biases and SSF's shift fold into the next bias, not into $W$. Two-sided scalings
    $\diag(a)W_0\diag(b)=W_0\odot(ab\T)$ lie in $\Im\mathrm{HiRA}_2$. As a morphism, $h(\ell)=(\ell-\mathbf 1,\mathbf
    1\T)$ lands in $\mathrm{HiRA}_1$ pointed at $(0,\mathbf 1\T)$, which is not HiRA's random-$A$ default pointing.
  \item \textbf{Orbit dimension.} The two-sided torus orbit $\{\diag(a)W_0\diag(b)\}$ has dimension $m+n-c(W_0)$, where
    $c(W_0)$ is the number of connected components of the bipartite support graph of $W_0$, an isolated vertex counting
    as its own component.
  \item \textbf{RoAd.} RoAd$_1$ (the paper's main variant, HF \texttt{road\_1}) is $W=RW_0$ with $R$ block-diagonal in
    scaled rotations $\alpha_i\mathrm{Rot}(\theta_i)$ of output pairs. It is the left action of
    $\CO^+(2)^{m/2}\cong(\mathbb C^\times)^{m/2}$: it is \textbf{(IA)$^{\mathbf 3}$ over $\mathbb C$}. It multiplies the
    input Gram matrix of each output pair by $\alpha_i^2$. Since $\alpha_i$ may be $0$, its image is the orbit closure,
    a real-linear space of dimension exactly $m$ whenever no output-pair block of $W_0$ is zero. RoAd$_2$ and RoAd$_4$
    (HF \texttt{road\_2}, \texttt{road\_4}) let each block range over all of $M_2(\R)$. They realise the block-diagonal
    action of the monoid $M_2(\R)^{m/2}$ (its group part $\GL_2^{m/2}$ is non-abelian), with linear image
    $\{\mathrm{blockdiag}(R_i)W_0\}$ of dimension $2m$ when every pair block has rank $2$. For them the Gram statement
    and ``over $\mathbb C$'' fail, and RoAd$_4$ carries a positive-dimensional gauge (two dimensions per row). HF pairs
    output $j$ with $j+g/2$ inside each group of size $g$ (default $64$) rather than adjacent outputs as in the paper,
    and also rotates the bias; the statements above do not depend on the pairing.
  \item \textbf{Zeros and covariance.} Every $W$ reachable by HiRA, (IA)$^3$ or SSF vanishes wherever $W_0$ does,
    $Z(W_0)\subseteq Z(W)$, and this persists under merging. New zeros can appear ($\ell_i=0$, or $(J+BA)_{ij}=0$), so
    the zero pattern itself is not invariant. The image covariance group of output-scaled boxes ((IA)$^3$ on $k,v$;
    post-linear SSF) is $\mathrm{Mon}_m\times\GL_n$, and that of input-scaled boxes ((IA)$^3$ on $W_{\mathrm{down}}$) is
    $\GL_m\times\mathrm{Mon}_n$. HiRA's is exactly $\mathrm{Mon}_m\times\mathrm{Mon}_n$ ($=\GL_1\times\GL_1$ when
    $m=n=1$).
\end{enumerate}
\end{proposition}

\paragraph{Numerical check}
The RoAd$_1$ orbit has dimension $8$ and the RoAd$_2$ image dimension $16$ at $m=8$ (Appendix~\ref{app:repro}).

\paragraph{Credit}
(IA)$^3$ is due to \citet{liu2022few}, SSF to \citet{lian2022scaling}, HiRA to \citet{huang2025hira}, RoAd to
\citet{liao20243}, and LayerNorm tuning to \citet{zhao2023tuning}. Feature-wise scaling as a building block goes back at
least to FiLM \citep{perez2017film}.

\begin{explains}
Each scaling method here is a rung $W_0\odot(J+Z)$: (IA)$^3$, SSF's scales and the gain part of LN-tuning on norms that
feed only linear maps lie in $\mathrm{HiRA}_1$, on one side or the other, and HiRA is LoRA in the Hadamard frame of
$W_0$. DoRA's image is the $\mathrm{Diag}_m$-orbit of LoRA's image (image level only; \thmref{Proposition}{II.8}). DoRA
and RoAd can fill zeros of $W_0$, so they become rungs only when $W_0$ is zero-free. Then every image germ can be written
as $W_0\odot(J+Z)$, so the ladder by itself constrains nothing there and serves as bookkeeping. RoAd$_1$ is complex
(IA)$^3$ and RoAd$_{2,4}$ are block-$M_2(\R)$ (IA)$^3$; for generic $W_0$ neither lies on a low-rank rung. HiRA's ``high
rank'' is a property of a frame: its transport $D_{W_0}$ depends on $\theta_0$ and does not preserve rank. Two re-based
cycles reach $W_0\odot(J+Z_1)\odot(J+Z_2)$, whose relative update $Z_1+Z_2+Z_1\odot Z_2$ has rank up to
$\min(2r+r^2,m,n)$, against $2r$ for LoRA (\thmref{Proposition}{II.13}).
\end{explains}

With HF defaults, (IA)$^3$ scales the output side of the attention projections it targets ($k$ and $v$, or $q$ and $v$)
and the input side of down, so the residual rotation above, acting on the other side, leaves its reachable functions
unchanged. HiRA's covariance group stops at $\mathrm{Mon}_m\times\mathrm{Mon}_n$ and fails the sufficient condition of
\thmref{Observation}{II.15}. That alone does not prove that its reachable functions change; Conjecture~A (stated after
\thmref{Observation}{II.15}) says they do, and \thmref{Prediction}{X.4} is the test.

Corollary~II.6, Theorem~II.7 and Propositions~II.8 and~II.9 were repaired after two rounds of adversarial review; the
referee record in \S\ref{sec:analogy} lists what changed.

The groups of this exposé see orbits. LoRA's cone is governed instead by the rank of the update, read off a cut in a
tensor network, while the size of an image is a dimension, which behaves like a volume. Rank and volume can move in
opposite directions, the subject of \S\ref{sec:cut}.
